\pagenumbering{arabic}
\setcounter{page}{1}

\chapter{Giới thiệu đề tài}

\section{Tổng quan tình hình liên quan đến đề tài}

Doanh nghiệp nhỏ và vừa (DNNVV) tại Việt Nam phải thực hiện đồng thời nhiều
nghĩa vụ pháp lý: thuế, lao động, bảo hiểm xã hội, kế toán và hợp đồng thương
mại. Cổng thông tin văn bản quy phạm pháp luật của nhà nước cung cấp nguồn gốc
đầy đủ và có tính pháp lý, nhưng cách tra cứu giả định người dùng đã biết số
hiệu văn bản hoặc thuật ngữ chuyên môn. Trên thực tế, kế toán và chủ doanh nghiệp
thường diễn đạt nhu cầu bằng ngôn ngữ đời thường, ví dụ ``thử việc được bao lâu''
hay ``quý này nộp GTGT ngày nào''. Kết quả tìm theo từ khóa dễ bị quá rộng hoặc
trượt mục tiêu, vì văn bản luật dùng cụm từ chuẩn hóa (``đơn phương chấm dứt hợp
đồng lao động'') khác với cách nói thông dụng (``sa thải'').

Các hệ thống hội thoại dựa trên mô hình ngôn ngữ lớn trả lời tiếng Việt trôi chảy
nhưng không bảo đảm dẫn chiếu đúng Điều. Khi đối chiếu với văn bản gốc, số Điều
có thể không tồn tại hoặc nội dung bị trộn từ văn bản khác. Với nghĩa vụ kèm chế
tài hành chính, câu trả lời sai nhưng mang tính khẳng định tạo rủi ro cao hơn so
với việc không có câu trả lời.

Về kỹ thuật, kiến trúc truy hồi-rồi-sinh (RAG)~\cite{lewis2020rag} buộc mô hình
sinh câu trên ngữ cảnh đã truy hồi. Dense retrieval~\cite{karpukhin2020dpr},
BM25~\cite{robertson2009bm25}, Reciprocal Rank Fusion~\cite{cormack2009rrf} và
HyDE~\cite{gao2023hyde} là các thành phần phổ biến để ghép pipeline hỏi đáp trên
corpus hẹp. Công trình chatbot tư vấn tuyển sinh dùng RAG trên tiếng
Việt~\cite{duong2026} cho thấy hướng này khả thi khi kho tri thức có biên rõ.
Đồ án áp dụng các kỹ thuật trên vào kho VBQPPL phục vụ DNNVV, với đơn vị văn bản
là Điều.

\section{Bối cảnh pháp lý và đặc thù DNNVV}

\subsection{Doanh nghiệp nhỏ và vừa trong hệ thống pháp luật Việt Nam}

Luật Hỗ trợ doanh nghiệp nhỏ và vừa số 04/2017/QH14 xác định DNNVV theo tiêu
chí quy mô (lao động, doanh thu, nguồn vốn) gắn với ngành nghề, và gắn một số
chính sách hỗ trợ~\cite{quochoi2017}. Đồ án không tái lập toàn bộ tiêu chí phân
loại để cấp ``giấy chứng nhận DNNVV''. Hồ sơ tổ chức trong hệ thống chỉ lưu các
tham số đủ để sinh lịch và mặc định chủ thể khi hỏi đáp: tên, mã số thuế, số
lao động, kỳ kê khai GTGT.

Điểm thực tiễn: nghĩa vụ thuế, bảo hiểm, báo cáo lao động và hóa đơn không biến
mất khi quy mô nhỏ. Bộ phận pháp chế riêng thì thường không có. Người dùng thật
của đồ án vì thế là chủ doanh nghiệp, kế toán dịch vụ hoặc nhân sự kiêm nhiệm
--- đối tượng hỏi bằng lời nói hàng ngày hơn là bằng số hiệu thông tư.

\subsection{Hệ thống văn bản quy phạm và thứ bậc áp dụng}

Văn bản do Quốc hội ban hành (luật, bộ luật) đặt nguyên tắc. Nghị định của
Chính phủ quy định chi tiết. Thông tư của bộ hướng dẫn thực hiện. Câu hỏi thao
tác (``nộp tờ khai ngày nào'', ``hóa đơn gửi khi nào'') thường rơi xuống nghị
định và thông tư, trong khi câu hỏi nguyên tắc (``có được đơn phương chấm dứt
không'') rơi ở luật. Corpus đồ án cố ý chứa cả ba cấp trong các lĩnh vực đã chọn,
tránh tình trạng trả lời đúng luật khung nhưng sai bước kê khai.

Hiệu lực theo thời gian là bài toán pháp điển riêng: văn bản mới thay văn bản
cũ, điều khoản chuyển tiếp, văn bản hết hiệu lực còn nằm trên cổng. Đồ án lưu
theo số hiệu, chưa gắn lịch sử hiệu lực từng Điều. Khi kho có hai luật cùng tên
khác số (ví dụ quản lý thuế, thuế GTGT), cả hai đều tìm được; người dùng và
người quản trị phải tự chọn bản đang áp dụng. Hạn chế này được nêu lại ở
Chương~5, không giấu trong phần ``đã hoàn thiện''.

\subsection{Rủi ro tuân thủ đối với người không chuyên}

Chậm nộp, sai tờ khai, thiếu báo cáo lao động có thể kéo theo tiền chậm nộp và
xử phạt hành chính. Chatbot không nguồn tạo rủi ro kiểu khác: câu tiếng Việt
trôi, kèm số Điều nghe quen, khiến người dùng tin hơn kết quả tìm từ khóa thô.
Đồ án đặt yêu cầu tối thiểu: mọi căn cứ đưa ra UI phải mở được về một Điều trong
kho. Yêu cầu này không biến hệ thống thành tư vấn pháp lý có trách nhiệm nghề
nghiệp. Người dùng vẫn phải đối chiếu văn bản gốc và, khi tranh chấp, nhờ tổ
chức hành nghề.

\subsection{Ranh giới trợ lý AI và dịch vụ pháp lý}

Hệ thống hỗ trợ tra cứu, nhắc việc và chỉ ra chỗ hợp đồng có thể lệch so với
Điều đã truy hồi. Hệ thống không đại diện tố tụng, không soạn hồ sơ nộp cơ quan
thuế, không cam kết ``thắng kiện''. Câu trả lời mang tính hỗ trợ thông tin trên
corpus đóng. Ranh giới này vừa là yêu cầu đạo đức nghề, vừa là phạm vi đồ án
ngành: thời gian không đủ để xây phân hệ chữ ký số hay kênh nộp thật.

\section{Bài toán tin học đặt ra từ hiện trạng}

Từ hiện trạng trên, bài toán không phải ``cài chatbot''. Ba lớp việc phải làm
cùng lúc:

\begin{enumerate}
    \item \textbf{Lớp dữ liệu.} Chuẩn hóa \soLuat{} văn bản thành Điều có siêu
    dữ liệu (số hiệu, chương, viện dẫn), nhập PostgreSQL, lập chỉ mục không dấu,
    nhúng vector.
    \item \textbf{Lớp suy diễn hạn chế.} Pipeline RAG lai, hậu kiểm citation,
    lịch theo rule --- suy diễn nằm ở chỗ ghép nguồn, không ở chỗ mô hình tự
    đặt ra hạn nộp.
    \item \textbf{Lớp phần mềm.} Đa người dùng, JWT, cách ly tổ chức, SSE, Docker.
    Không có lớp này thì pipeline chỉ chạy được trên máy phát triển.
\end{enumerate}

Báo cáo lần lượt trả lời ba lớp: Chương~2 cho kỹ thuật truy hồi và công nghệ;
Chương~3 cho phân tích thiết kế; Chương~4 cho bằng chứng chạy được.


Sản phẩm là ứng dụng web đa người dùng: đăng nhập, lịch sử hội thoại, tra cứu
văn bản, hỏi đáp có dẫn chiếu, soát xét hợp đồng và lịch tuân thủ. Pipeline hỏi
đáp cắt corpus theo Điều, truy hồi lai BM25--vector, rồi sinh câu trên ngữ cảnh
đã lọc. Mô hình mặc định là Gemini (lớp tương thích OpenAI). Đánh giá đồ án dựa
trên kịch bản vận hành Docker và khả năng đối chiếu Điều gốc, không dựa trên
bảng xếp hạng mô hình bên ngoài.

\section{Lý do lựa chọn đề tài}

DNNVV thường không duy trì bộ phận pháp chế riêng, trong khi nghĩa vụ định kỳ
không giảm theo quy mô: kê khai thuế GTGT, tạm nộp và quyết toán TNDN, đóng BHXH,
báo cáo lao động, nộp báo cáo tài chính năm. Các mốc này lặp lại, có hạn chót và
gắn với chế tài khi chậm. Một công cụ hỗ trợ hỏi đáp có dẫn chiếu, tra cứu Điều
gốc, soát hợp đồng và nhắc lịch nộp phù hợp hơn so với chatbot chỉ đưa lời khuyên
chung.

Phạm vi được giới hạn ở thuế, lao động, bảo hiểm, kế toán, hợp đồng dân sự --
thương mại và sở hữu trí tuệ mức cơ bản. Tố tụng, hình sự và đất đai chuyên sâu
không thuộc đề tài, nhằm bảo đảm hiện thực và kiểm chứng được trong thời gian đồ
án ngành.

Về kỹ thuật, ba bài toán không giải được bằng cách cài một chatbot mẫu:

\begin{enumerate}
    \item \textbf{Lệch từ vựng giữa người dùng và văn bản luật.} Truy hồi chỉ
    theo từ khóa dễ bỏ sót Điều đúng; truy hồi chỉ theo vector dễ nhầm Điều kề
    nhau. Cần kết hợp hai hướng rồi gộp hạng.
    \item \textbf{Tìm kiếm tiếng Việt trên PostgreSQL.} Hệ quản trị này không có
    từ điển full-text tiếng Việt chuẩn. Người dùng thường gõ không dấu. Nếu tìm
    và hiển thị trên cùng bản đã bỏ dấu, đoạn trích mất dấu, khó đọc. Cần tách
    pha tìm (bản không dấu) và pha hiển thị (bản còn dấu).
    \item \textbf{Hiện tượng bịa số Điều.} Nếu lấy đầu ra mô hình làm nguồn
    chính, hệ thống không khác chatbot không kiểm chứng. Cần hậu kiểm: số Điều
    trong câu trả lời phải thuộc tập vừa truy hồi.
\end{enumerate}

\section{Mục tiêu đề tài}

\subsection{Mục tiêu tổng quát}

Xây dựng ứng dụng web đa người dùng, đa tổ chức, hỗ trợ DNNVV tra cứu và hỏi đáp
pháp lý, trong đó mọi căn cứ pháp lý hiển thị đều truy được về Điều trong corpus
đã chuẩn hóa.

\subsection{Mục tiêu cụ thể}

\begin{enumerate}
    \item Xây dựng corpus DNNVV gồm \soLuat{} văn bản, \soDieu{} Điều,
    \soChuong{} chương; cắt theo Điều; lưu metadata và viện dẫn chéo.
    \item Hiện thực pipeline hỏi đáp bảy bước: phân tích ý định, chuẩn bị truy
    vấn (HyDE), truy hồi lai, lọc ngữ cảnh, sinh câu, hậu kiểm trích dẫn; truyền
    kết quả theo SSE.
    \item Xây dựng chức năng tra cứu theo cây Chương--Điều, tìm không dấu, hỗ trợ
    hỏi tiếp từ Điều đang đọc.
    \item Xây dựng chức năng soát xét hợp đồng PDF/DOCX/TXT, xử lý nền, theo dõi
    trạng thái, xuất bảng rủi ro.
    \item Xây dựng lịch tuân thủ \soRule{} loại nghĩa vụ, sinh theo hồ sơ doanh
    nghiệp, kèm dẫn chiếu pháp lý.
    \item Hiện thực xác thực JWT (access và refresh), bốn vai trò, cách ly dữ
    liệu theo người dùng và theo tổ chức.
    \item Đóng gói bằng Docker Compose (ba dịch vụ).
\end{enumerate}

\section{Cách tiếp cận và phương pháp}

\subsection{Hỏi đáp dạng open-book}

Hỏi đáp close-book dựa vào tri thức nén trong tham số mô hình, không chỉ được
văn bản nguồn và dễ sai khi luật thay đổi. Đồ án chọn open-book: mọi kết luận
pháp lý phải dựa trên đoạn văn đã truy hồi từ corpus~\cite{lewis2020rag}. Khi
không tìm được Điều liên quan, hệ thống thông báo thiếu căn cứ, không sinh đoạn
khuyên chung giả danh quy định.

\subsection{Truy hồi lai}

Truy hồi dày (vector) bắt ngữ nghĩa; BM25 bắt mặt chữ và số hiệu. Reciprocal Rank
Fusion gộp hai bảng xếp hạng mà không cần đưa điểm về cùng thang
đo~\cite{cormack2009rrf}. Cấu hình đồ án: trọng số dense:BM25 $= 3:1$,
$k_{\mathrm{RRF}}=60$, lấy $k=8$ Điều cho các bước sau.

\subsection{Quy trình phát triển}

Hệ thống được phát triển theo vòng ngắn: chuẩn hóa corpus, xây dựng CSDL và bảng
tri thức, hiện thực xác thực và hội thoại, lần lượt hoàn thiện tra cứu, hợp đồng
và lịch tuân thủ. Mỗi module hoàn thành API được gắn giao diện ngay, nhằm phát
hiện sớm các lỗi chỉ xuất hiện khi vận hành (định tuyến số hiệu chứa dấu
\texttt{/}, thứ tự bản ghi cùng mốc thời gian, tải tệp multipart).

Tìm kiếm toàn văn phụ thuộc \texttt{tsvector}, \texttt{unaccent} và
\texttt{pg\_trgm} trên PostgreSQL, không thay bằng SQLite.

\section{Đối tượng và phạm vi}

\subsection{Đối tượng}

\begin{itemize}
    \item Đối tượng sử dụng: chủ DNNVV, kế toán, nhân sự --- tra cứu và hỏi đáp
    bằng tiếng Việt thông thường.
    \item Đối tượng kỹ thuật: pipeline RAG lai trên corpus VBQPPL cắt theo Điều;
    ứng dụng web đa người dùng.
\end{itemize}

\subsection{Phạm vi chức năng}

Bốn nghiệp vụ: hỏi đáp có trích dẫn, tra cứu cây văn bản, soát xét hợp đồng, lịch
tuân thủ. Đồ án không thực hiện OCR tệp ảnh, không nộp hồ sơ thuế hộ, không ký số
và không đại diện tố tụng. Lịch tuân thủ là công cụ nhắc việc, không thay thế phần
mềm kế toán.

\subsection{Phạm vi dữ liệu và kỹ thuật}

Corpus gồm \soLuat{} văn bản thuộc thuế, lao động, BHXH, kế toán, hợp đồng, sở hữu
trí tuệ mức cơ bản và hỗ trợ DNNVV. Luật hình sự và văn bản địa phương không nằm
trong kho. Mô hình ngôn ngữ và embedding được gọi qua mạng; đồ án không tự triển
khai GPU. Soát xét hợp đồng dùng hàng đợi nội bộ \texttt{BackgroundTasks} của
FastAPI, phù hợp quy mô đồ án, chưa phải hàng đợi phân tán.

\section{Tổng quan các giải pháp hiện có}

\begin{table}[htp]
\centering
\caption{Đối chiếu cổng VBQPPL, chatbot không nguồn và hướng đồ án.}
\label{tab:khao-sat}
\begin{tabular}{p{2.8cm}p{3.6cm}p{3.6cm}p{3.8cm}}
\toprule
\textbf{Tiêu chí} & \textbf{Cổng VBQPPL} & \textbf{Chatbot không nguồn} &
\textbf{Hướng đồ án} \\
\midrule
Cách hỏi & Số hiệu, từ khóa & Câu nói thường & Câu thường hoặc từ cây văn bản \\
Đối chiếu nguồn & Chính là văn bản luật & Không có dẫn chiếu mở được & Mỗi căn cứ
mở đúng Điều \\
Bịa số Điều & Không xảy ra & Dễ gặp & Loại căn cứ ngoài tập truy hồi \\
Nghiệp vụ kế toán, nhân sự & Không & Không & Soát hợp đồng, lịch nộp \\
Độ rộng kho & Rất rộng & Không xác định nguồn & \soLuat{} văn bản, phạm vi hẹp \\
\bottomrule
\end{tabular}
\end{table}

Cổng VBQPPL của nhà nước vượt trội về độ tin cậy và độ phủ. Chatbot ngôn ngữ lớn
vượt trội về cách hỏi tự nhiên. Đồ án không cạnh tranh độ phủ với cổng nhà nước.
Hướng tiếp cận là kết hợp câu hỏi ngôn ngữ tự nhiên với kho văn bản hẹp nhưng mở
được từng Điều, đồng thời bổ sung hai nghiệp vụ mà cổng tra cứu không cung cấp:
soát xét hợp đồng và nhắc mốc tuân thủ.

\section{Ý nghĩa của đề tài}

Đồ án không đề xuất thuật toán truy hồi mới. Ý nghĩa thực tiễn nằm ở việc chuyển
pipeline hỏi đáp thành hệ thống vận hành được cho DNNVV, cụ thể:

\begin{itemize}
    \item Ứng dụng đa người dùng, đa tổ chức, có xác thực và cách ly dữ liệu,
    thay cho kịch bản hỏi đáp chạy một lần.
    \item Tra cứu tiếng Việt không dấu trên PostgreSQL trong khi đoạn trích hiển
    thị vẫn còn dấu; lịch tuân thủ sinh theo hồ sơ từng doanh nghiệp và dẫn được
    Điều có trong corpus.
    \item Hệ thống suy giảm có kiểm soát: khi thiếu khóa mô hình, tra cứu và lịch
    tuân thủ vẫn dùng được; hỏi đáp và soát xét hợp đồng báo chưa sẵn sàng.
\end{itemize}

\section{Phương pháp nghiên cứu và công cụ}

Đồ án kết hợp nghiên cứu tài liệu và hiện thực hệ thống. Phần tài liệu gồm các
công trình RAG, BM25, dense retrieval, RRF, HyDE và giáo trình công nghệ phần
mềm~\cite{thanh2024}. Phần hiện thực tuân thủ vòng đời: phân tích yêu cầu, thiết
kế kiến trúc, cài đặt, vận hành thử, đóng gói.

Công cụ chính: Visual Studio Code / Cursor để soạn mã; Git để lưu phiên bản;
Docker Compose để dựng môi trường; XeLaTeX để biên soạn báo cáo. Dữ liệu pháp lý
lấy từ văn bản đã công bố, không thu thập dữ liệu cá nhân ngoài tài khoản thử
nghiệm.

\section{Bố cục báo cáo}

Báo cáo gồm phần đầu (lời cảm ơn, nhận xét, tóm tắt tiếng Việt, abstract tiếng
Anh, mục lục, danh mục), năm chương, tài liệu tham khảo và phụ lục.

Chương~1 đặt bài toán DNNVV, mục tiêu, phạm vi và đối chiếu giải pháp hiện có.
Chương~2 trình bày lý thuyết và công nghệ đồ án thực sự dùng. Mỗi mục đi theo
bốn câu (là gì, vì sao, cách dùng, khi nào) và lấy ví dụ pháp lý tiếng Việt,
không chép định nghĩa bách khoa. Chương~3 phân tích yêu cầu,
use case, thiết kế CSDL, API, danh mục \soLuat{} văn bản và \soRule{} rule tuân
thủ. Chương~4 mô tả môi trường thử, kịch bản Docker và tiêu chí định tính.
Chương~5 kết luận, hạn chế, hướng phát triển.

Phụ lục hướng dẫn chạy hệ thống, thuật ngữ ngắn, quy trình nạp corpus và ràng
buộc sử dụng.
